\section{总结与展望}
\label{sec:conclusion}
NNPDF3.1 是 NNPDF 家族新的主版本 PDF 发布。
%
通过纳入许多新观测量的约束，它相对于 NNPDF3.0 是一大进步；得益于相应 NNLO QCD 修正在近期变为可用，
其中若干观测量首次进入全局 PDF 确定。突出的例子有 $t\bar{t}$ 微分分布与 $Z$ 玻色子 $p_T$ 谱。
%
从理论角度看，最主要的改进是把粲 PDF 放到与轻夸克 PDF 同等的地位上。
%
独立参数化粲 PDF 化解了 ATLAS 规范玻色子产生与 HERA 单举结构函数数据之间本会存在的张力，
改善了与 LHC 数据的符合程度，并把微扰生成粲对极点粲质量的强依赖转化为 PDF 不确定度——因为大部分质量依赖被重新吸收进初始 PDF 形状。

NNPDF3.1 集合也是首个以多种格式交付 PDF 的集合：首先它们同时以 Hessian 与蒙特卡罗形式发布；
其次默认集合经过优化与压缩，使得较少的蒙特卡罗副本或 Hessian 误差集合即可重现大得多的底层副本集合的统计特征。
下面我们讨论如何从大量蒙特卡罗副本产生 Hessian 与蒙特卡罗缩减集合；
随后总结通过 LHAPDF 接口公开的全部 PDF 集合；
最后对未来发展作简短展望。

\subsection{NNPDF3.1 缩减集合的验证}
\label{sec:compression}

  默认 NNPDF3.1 NLO 与 NNLO PDF（中心值 $\alpha_s(m_Z)=0.118$）连同第~\ref{sec:results-mc}
  节讨论的微扰粲修改版，都以 $N_{\rm rep}=1000$ 副本集合产生。这些大副本样本随后经两种缩减策略处理：
  其一是压缩蒙特卡罗（CMC）算法~\cite{Carrazza:2015hva}，得到基于较少副本数的蒙特卡罗表示；
  其二是 MC2H 算法~\cite{Carrazza:2015aoa}，以固定数目的误差集合实现底层 PDF 概率分布的最优 Hessian 表示。
  具体地，我们由此构造了 $N_{\rm rep}=100$ 副本的 CMC-PDF 集合与 $N_{\rm eig}=100$ 个（对称）本征矢量的 MC2H 集合。

  图~\ref{fig:cmc-mc2h-validation} 比较 NNPDF3.1 NNLO 输入集合（$N_{\rm rep}=1000$ 副本）的 PDF
        与相应的缩减集合：$N_{\rm rep}=100$ 副本的 CMC-PDF 与 $N_{\rm eig}=100$
        本征值的 MC2H Hessian PDF。
        %
        各种情形下输入的 $N_{\rm rep}=1000$ 副本 MC PDF 与两个缩减集合都符合很好。
        %
        由构造可知，MC2H 集合在中心值与 PDF 方差上的符合略好，
        因为 CMC-PDF 集合还力图重现概率分布的更高矩乃至可能的非高斯特征，而 Hessian 集合按构造是高斯的。
        %
        沿用文献~\cite{Carrazza:2015aoa,Carrazza:2015hva}的分析，
        我们核实了 PDF 之间的关联也被高度准确地重现。
        
%%%%%%%%%%%%%%%%%%%%%%%%%%%%%%%%%%%%%%%%%%%%%%%%%%%%%%%%%%%%%%%%%%%%%
\begin{figure}[t]
  \begin{center}
     \includegraphics[scale=0.32]{images/plots/xu-cmc-mc2h-nnlo-fc.pdf}
     \includegraphics[scale=0.32]{images/plots/xd-cmc-mc2h-nnlo-fc.pdf}
     \includegraphics[scale=0.32]{images/plots/xubar-cmc-mc2h-nnlo-fc.pdf}
     \includegraphics[scale=0.32]{images/plots/xdbar-cmc-mc2h-nnlo-fc.pdf}
     \includegraphics[scale=0.32]{images/plots/xsp-cmc-mc2h-nnlo-fc.pdf}
     \includegraphics[scale=0.32]{images/plots/xc-cmc-mc2h-nnlo-fc.pdf}
     \includegraphics[scale=0.32]{images/plots/xsinglet-cmc-mc2h-nnlo-fc.pdf}
     \includegraphics[scale=0.32]{images/plots/xg-cmc-mc2h-nnlo-fc.pdf}
     \caption{\small 三者的比较：$N_{\rm rep}=1000$ 副本的 NNPDF3.1 NNLO 输入集合 PDF、
       $N_{\rm rep}=100$ 副本的压缩蒙特卡罗 CMC-PDF、以及 $N_{\rm eig}=100$ 对称
       本征值的 MC2H Hessian PDF。
    \label{fig:cmc-mc2h-validation}}
\end{center}
\end{figure}
%%%%%%%%%%%%%%%%%%%%%%%%%%%%%%%%%%%%%%%%%%%%%%%%%%%%%%%%%%%%%%%%%%%%%%%

为了验证 CMC-PDF 算法从初始 $N_{\rm rep}=1000$ 副本压缩到 $N_{\rm rep}=100$ 副本的效率，
图~\ref{fig:cmc-mc2h-validation-nrep} 按照文献~\cite{Carrazza:2015hva}描述的程序给出统计估计量的汇总：
比较由 $N_{\rm rep}=1000$ 副本的 NNPDF3.1 NNLO 输入集合定义的概率分布的特定性质与相应压缩集合的性质，
表示为缩减集合副本数 $\widetilde{N}_{\rm rep}$（从 $\widetilde{N}_{\rm rep}=100$ 起）的函数。
        %
        我们把压缩算法的结果与从原 1000 个副本随机抽取 $\widetilde{N}_{\rm rep}$ 个的做法比较：
        图中给出对应中心值、标准差、峰度与偏度、关联以及 Kolmogorov 距离的误差函数 ERF。
        %
        这些结果表明一个 100 副本的 CMC-PDF 集合大致重现了随机抽取的 $\widetilde{N}_{\rm rep}=400$
        PDF 集合所含的信息。


%%%%%%%%%%%%%%%%%%%%%%%%%%%%%%%%%%%%%%%%%%%%%%%%%%%%%%%%%%%%%%%%%%%%%
\begin{figure}[t]
  \begin{center}
    \includegraphics[scale=0.35]{images/plots/cmc_cv.pdf}\includegraphics[scale=0.35]{images/plots/cmc_std.pdf}
    \includegraphics[scale=0.35]{images/plots/cmc_skewness.pdf}\includegraphics[scale=0.35]{images/plots/cmc_kurtosis.pdf}
    \includegraphics[scale=0.35]{images/plots/cmc_kolmogorov.pdf}\includegraphics[scale=0.35]{images/plots/cmc_correlation.pdf}

     \caption{\small 由 $N_{\rm
          rep}=1000$ 副本的 NNPDF3.1 NNLO 输入集合计算的与由 $\widetilde{N}_{\rm rep}$ 副本压缩集合计算的
       概率分布估计量的比较，表示为 $\widetilde{N}_{\rm rep}$ 的函数。
       图示为对应中心值、标准差、峰度、偏度、关联与 Kolmogorov 距离的误差函数（ERF）。
    \label{fig:cmc-mc2h-validation-nrep}
  }
\end{center}
\end{figure}
%%%%%%%%%%%%%%%%%%%%%%%%%%%%%%%%%%%%%%%%%%%%%%%%%%%%%%%%%%%%%%%%%%%%%%%

\clearpage



\subsection{交付清单}
\label{sec:delivery}

现在完整列出将通过 {\sc\small LHAPDF6}
接口~\cite{Buckley:2014ana}公开发布的 NNPDF3.1 PDF 集合：
\begin{center}
{\bf \url{http://lhapdf.hepforge.org/}~} 。
\end{center}
正如文中反复提到的，与这些 PDF 集合相关的大量结果可从存储库获得：
\begin{center}
{\bf \url{http://nnpdf.hepforge.org/html/nnpdf31/catalog}~} 。
\end{center}
所有集合都以 $N_{\rm rep}=100$ 的蒙特卡罗集合提供。
对基线集合而言，它们由更大的 $N_{\rm rep}=1000$ 副本集合构造而来，后者同样公开。
基线集合还以含 100 个误差集合的 Hessian 集合形式提供。

完整清单如下：

\begin{itemize}

\item {\it 基线 NLO 与 NNLO NNPDF3.1 集合}

基线 NLO 与 NNLO NNPDF3.1 集合基于全局数据集，$\alpha_s(m_Z)=0.118$，
变味数方案最多五个活跃味道。这些集合含 $N_{\rm rep}=1000$ 个 PDF 副本。
\begin{flushleft}
\tt NNPDF31\_nlo\_as\_0118\_1000 \\
\tt NNPDF31\_nnlo\_as\_0118\_1000.\\
\end{flushleft}
粲为微扰生成（如同以往 NNPDF 集合）的修改版也已发布，同样含 $N_{\rm rep}=1000$ 个 PDF 副本：
\begin{flushleft}
\tt NNPDF31\_nlo\_pch\_as\_0118\_1000 \\
\tt NNPDF31\_nnlo\_pch\_as\_0118\_1000.\\
\end{flushleft}
由它们出发构造了优化的 $N_{\rm rep}=100$ 蒙特卡罗集合
\begin{flushleft}
\tt NNPDF31\_nlo\_as\_0118 \\
\tt NNPDF31\_nnlo\_as\_0118 \\
\tt NNPDF31\_nlo\_pch\_as\_0118 \\
\tt NNPDF31\_nnlo\_pch\_as\_0118,
\end{flushleft}
以及 $N_{\rm eig}=100$ 本征矢量的 Hessian 集合
\begin{flushleft}
\tt NNPDF31\_nlo\_as\_0118\_hessian \\
\tt NNPDF31\_nnlo\_as\_0118\_hessian \\
\tt NNPDF31\_nlo\_pch\_as\_0118\_hessian \\
\tt NNPDF31\_nnlo\_pch\_as\_0118\_hessian,
\end{flushleft}
如上文第~\ref{sec:compression} 节所述。

在这些基础上，可用 {\tt SM-PDF} 工具~\cite{Carrazza:2016htc}
构造为特定观测量计算而优化的小型本征矢量集合。
具体地，可为一系列预定义观测量生成优化集合并经由网页界面~\cite{Carrazza:2016wte}下载：
\begin{center}
{\bf \url{https://smpdf.mi.infn.it}~} 。
\end{center}





\item {\it 味数变化}

我们产生了 NLO 与 NNLO 的集合，其最大味数不同于默认的 5：或扩展至六，或冻结在四：
\begin{flushleft}
  \tt NNPDF31\_nlo\_as\_0118\_nf\_4 \\
  \tt NNPDF31\_nlo\_as\_0118\_nf\_6 \\
  \tt NNPDF31\_nnlo\_as\_0118\_nf\_4 \\
  \tt NNPDF31\_nnlo\_as\_0118\_nf\_6. \\
\end{flushleft}
微扰生成粲变体同样可用，此时还有 $n_f=3$ 固定味数方案：
\begin{flushleft}
  \tt NNPDF31\_nlo\_pch\_as\_0118\_nf\_3 \\
  \tt NNPDF31\_nlo\_pch\_as\_0118\_nf\_4 \\
  \tt NNPDF31\_nlo\_pch\_as\_0118\_nf\_6 \\
  \tt NNPDF31\_nnlo\_pch\_as\_0118\_nf\_3 \\
  \tt NNPDF31\_nnlo\_pch\_as\_0118\_nf\_4 \\
  \tt NNPDF31\_nnlo\_pch\_as\_0118\_nf\_6. \\
\end{flushleft}




\item {\it $\alpha_s$ 变化。}

我们产生了下列 $\alpha_s(m_Z)$ 值的 NNLO 集合：0.108、0.110、0.112、0.114、0.116、0.117、0.118、0.119、0.120、0.122、0.124（原文首项"1.110"系笔误）：
\begin{flushleft}
  \tt NNPDF31\_nnlo\_as\_0108 \\
  \tt NNPDF31\_nnlo\_as\_0110 \\
  \tt NNPDF31\_nnlo\_as\_0112 \\
  \tt NNPDF31\_nnlo\_as\_0114 \\
  \tt NNPDF31\_nnlo\_as\_0116 \\
  \tt NNPDF31\_nnlo\_as\_0117 \\
  \tt NNPDF31\_nnlo\_as\_0118 \\
  \tt NNPDF31\_nnlo\_as\_0119 \\
  \tt NNPDF31\_nnlo\_as\_0120 \\
  \tt NNPDF31\_nnlo\_as\_0122 \\
  \tt NNPDF31\_nnlo\_as\_0124. \\
\end{flushleft}

对 $\alpha_s(m_Z)=0.116$ 与 $\alpha_s(m_Z)=0.120$ 还产生了 NLO 集合：
\begin{flushleft}
  \tt NNPDF31\_nlo\_as\_0116 \\
  \tt NNPDF31\_nlo\_as\_0120, \\
%  \tt NNPDF31\_nnlo\_as\_0116 \\
%  \tt NNPDF31\_nnlo\_as\_0120,\\
\end{flushleft}
以及 NLO 与 NNLO 的微扰粲变体：
\begin{flushleft}
  \tt NNPDF31\_nlo\_pch\_as\_0116 \\
  \tt NNPDF31\_nlo\_pch\_as\_0120 \\
\tt NNPDF31\_nnlo\_pch\_as\_0116\\
\tt NNPDF31\_nnlo\_pch\_as\_0120.\\
\end{flushleft}


为便于计算 PDF+$\alpha_s$ 组合不确定度，我们还提供了针对 $\alpha_s(m_Z)=0.118\pm0.002$ 的
PDF+$\alpha_s$ 打包变化集合。它们既有蒙特卡罗形式也有 Hessian 形式。
\begin{flushleft}
  \tt
NNPDF31\_nlo\_pdfas	\\
NNPDF31\_nnlo\_pdfas	\\
NNPDF31\_nlo\_pch\_pdfas \\	
NNPDF31\_nnlo\_pch\_pdfas  \\	
NNPDF31\_nlo\_hessian\_pdfas \\
NNPDF31\_nnlo\_hessian\_pdfas \\	
NNPDF31\_nlo\_pch\_hessian\_pdfas \\	
NNPDF31\_nnlo\_pch\_hessian\_pdfas. \\
\end{flushleft}
它们的构造方式如下：
\begin{enumerate}
\item 中心值（PDF 成员 0）即对应 $\alpha_s(m_Z)=0.118$ 集合的中心值。
\item PDF 成员 1 至 100 对应来自 $\alpha_s(m_Z)=0.118$ 集合的 $N_{\rep}=100$
  （$N_{\rm eig}=100$）个蒙特卡罗副本（Hessian 本征矢量）。
\item PDF 成员 101 与 102 分别是 $\alpha_s(m_Z)=0.116$ 与 $\alpha_s(m_Z)=0.120$ 集合的中心值。
\end{enumerate}
注意因此在 Hessian 情形成员 0 是中心集合，其余打包成员 $1-102$ 都是误差集合；
而在蒙特卡罗情形，成员 $1-100$ 为蒙特卡罗副本，成员 0、$101$、$102$
为中心集合（副本平均）。如何用它们计算 PDF+$\alpha_s$ 组合不确定度见例如文献~\cite{Butterworth:2015oua}。


\item {\it 粲质量变化}


我们提供不同极点粲质量 $m_c^{\rm pole}$ 取值的集合。
  它们仅在 NNLO 提供，取 $m_c^{\rm pole}=1.38$ GeV 与 $m_c^{\rm pole}=1.64$ GeV。
  \begin{flushleft}
    {\tt NNPDF31\_nnlo\_as\_0118\_mc\_138}\\
    {\tt NNPDF31\_nnlo\_as\_0118\_mc\_164}.\\
\end{flushleft}
供比较，相应的微扰粲修改版也可用：
  \begin{flushleft}
     {\tt NNPDF31\_nnlo\_pch\_as\_0118\_mc\_138}\\
    {\tt NNPDF31\_nnlo\_pch\_as\_0118\_mc\_164}.\\
\end{flushleft}


\item {\it 强制正定性集合}

我们提供 PDF 非负的集合：
  \begin{flushleft}
    {\tt NNPDF31\_nlo\_as\_0118\_mc}\\
     {\tt NNPDF31\_nlo\_pch\_as\_0118\_mc}\\
     {\tt NNPDF31\_nnlo\_as\_0118\_mc}\\
      {\tt NNPDF31\_nnlo\_pch\_as\_0118\_mc}\\
\end{flushleft}
它们的构造很简单：凡 PDF 出现负值即置零。因此这只是一种近似，
为方便与那些遇到负 PDF 即失效的程序配合使用而提供。

\item {\it LO 集合} 

领头阶 PDF 集合以 $\alpha_s = 0.118$ 与 $\alpha_s = 0.130$ 提供：
  \begin{flushleft}
    {\tt NNPDF31\_lo\_as\_0118}\\
    {\tt NNPDF31\_lo\_as\_0130}\\
\end{flushleft}
相应的微扰粲变体同样提供：
  \begin{flushleft}
    {\tt NNPDF31\_lo\_pch\_as\_0118}\\
    {\tt NNPDF31\_lo\_pch\_as\_0130}.
\end{flushleft}


\item {\it 缩减数据集}

从完整 NNPDF3.1 子集确定的 PDF（第~\ref{sec:impactzpt}--\ref{sec:collideronly} 节讨论）同样可用，具体为：
  \begin{flushleft}
  \begin{tabular}{ll}
  {\tt NNPDF31\_nnlo\_as\_0118\_noZpt} & 无 $Z$ $p_T$ 数据，见第~\ref{sec:impactzpt} 节； \\
  {\tt NNPDF31\_nnlo\_as\_0118\_notop} & 无 $t\bar{t}$ 产生数据，见第~\ref{sec:impacttop} 节； \\
  {\tt NNPDF31\_nnlo\_as\_0118\_nojets} & 无喷注数据，见第~\ref{sec:jetdata} 节； \\
  {\tt NNPDF31\_nnlo\_as\_0118\_wEMC} & 加入 EMC 粲数据，见第~\ref{sec:emc} 节；\\
  {\tt NNPDF31\_nnlo\_as\_0118\_noLHC} & 无 LHC 数据，见第~\ref{sec:nolhc} 节；  \\
  {\tt NNPDF31\_nnlo\_as\_0118\_proton}  & 仅质子靶数据，见第~\ref{sec:nonucl} 节；\\
  {\tt NNPDF31\_nnlo\_as\_0118\_collider} & 仅对撞机数据，见第~\ref{sec:collideronly} 节。\\
  \end{tabular}
  \end{flushleft}
\end{itemize}

除此之外，本文讨论的其他任何 PDF 集合均可应请求提供。

\subsection{展望}
\label{sec:outlook}

这里给出的 NNPDF3.1 PDF 确定是对 NNPDF3.0 的更新，却在方法学与数据集两方面都包含实质创新，
产生了精度与准确性都显著提高的 PDF 集合。得益于此，可以推进多项衍生项目：
或旨在更新以往基于旧版 NNPDF 发布的结果，或因为更高的准确性使以前不可能或不感兴趣的项目变得可行。

这些衍生项目包括：
\begin{itemize}
\item 强耦合常数 $\alpha_s(m_Z)$ 的精密确定。
%
与既往 NNPDF 确定~\cite{Lionetti:2011pw,Ball:2011us}相比，我们预期精度与准确性都会显著提高。
具体地，得益于在 NNLO 纳入众多同时对胶子与 $\alpha_s(m_Z)$ 有直接依赖的对撞机观测量，
彻底弃用核靶与低标度数据可能反而有利，即把 $\alpha_s$ 确定建立在第~\ref{sec:collideronly} 节的
仅对撞机数据集、甚或更保守的数据集之上。
\item 粲质量 $m_c$ 的确定。这将首次基于一个粲 PDF 被独立参数化的 PDF 确定，
从而避免把 $m_c$ 认作决定粲 PDF 大小的参数所带来的偏差。
  %
\item 包含光子 PDF $\gamma(x,Q^2)$ 的 NNPDF3.1QED 集合，从而更新以往的
  NNPDF2.3QED~\cite{Ball:2013hta} 与 NNPDF3.0QED~\cite{Bertone:2016ume} 集合。
  该更新应纳入近期的理论进展：具体包括直接从核子结构函数确定光子 PDF 的“LuxQED”
  约束~\cite{Manohar:2016nzj}；
以及现已包含在 {\tt APFEL}~\cite{Bertone:2013vaa}中的 PDF 演化与 DIS 系数函数的 NLO QED 修正。
\item 包含小 $x$ 重求和的 PDF 集合，基于文献~\cite{Altarelli:2005ni,Ball:2007ra,Altarelli:2008aj}
发展的、已在 {\tt HELL} 程序~\cite{Bonvini:2016wki}中实现并联接 {\tt APFEL} 的形式。
%
这将回应一个长期悬而未决的问题：固定阶微扰演化能否恰当描述小 $x$ 的单举 HERA
数据~\cite{Caola:2010cy,Caola:2009iy}，并最终有助于更好地控制小 $x$ 理论不确定度。
\end{itemize}

更长远的时期里，预期数据集与方法学都将有进一步的重要改进。
一方面，到目前为止我们基本上局限于 LHC 8 TeV 数据。未来版本将纳入大量 13 TeV 测量——
其中若干来自 ATLAS、CMS 与 LHCb 的测量已经可用。具体设想纳入更多目前不在数据集之内、
对 PDF 有重大潜在影响且高阶修正已经具备的过程。其中包括瞬发光子产生~\cite{Aaboud:2017cbm}——
其 NNLO 修正现已成为可用~\cite{Campbell:2016lzl}，且其对 PDF 的影响众所周知~\cite{d'Enterria:2012yj}；
单顶产生（NNLO 也已知~\cite{Brucherseifer:2014ama}）；以及可能的前向 $D$ 介子产生或更一般的含末态
$D$ 介子的过程，例如最近由 ATLAS 测量的 $W+D$ 产生~\cite{Aad:2014xca}——
其对 PDF 的重要性曾被反复强调~\cite{Zenaiev:2015rfa,Gauld:2015yia,Gauld:2016kpd}。

数据集的进一步扩充很可能要求 PDF 确定方法学的实质性改进。同时也很可能导致 PDF 不确定度的进一步缩小，
从而要求对理论不确定度更好的控制。我们特别预期两个方向上的重要进展。
其一，目前尚未纳入、靠运动学切割加以控制的电弱修正，必须以更系统的方式纳入。
其二，缺失高阶修正导致的理论不确定度也必须予以估计。事实上第~\ref{sec:thunc} 节的初步估计表明：
当前未计入 PDF 不确定度的缺失高阶不确定度很可能很快变得不可忽略，在某些运动学区域甚至会占据主导。
所有这些改进都将成为未来一次重大 PDF 发布的组成部分。

\bigskip
\bigskip
\begin{center}
\rule{5cm}{.1pt}
\end{center}
\bigskip
\bigskip
